\section{问题重述}

\subsection{问题背景}

小区微网由负载、光伏发电、储能装置和外部电网共同构成。光伏出力具有明显的日内规律和随机偏差，外部电价既可能采用固定日内曲线，也可能随日期和时段实时波动。储能能够在低价或光伏富余时充电、在高价或供电不足时放电，但其容量、功率和转换效率均受限制。因此，调度策略必须同时处理预测误差、跨时段能量转移和高价紧急购电风险。

\subsection{问题内容}

\textbf{问题一：}根据固定的一日电价、负载和光伏预测，制定 144 个 10 分钟区间的计划购电量，使购电费最小，并保证日首与日末储电量相同。

\textbf{问题二：}面对全年变化的实际负载和光伏，在每天 0:00 仅利用历史信息制定计划；供电不足时允许按当时电价的 5 倍紧急购电，并给出 2 月 1 日至 12 月 31 日的完整策略。

\textbf{问题三：}利用每日 0:00、6:00、12:00、18:00 发布的未来 24 小时光伏预报滚动调整尚未执行的购电计划，计入上调、下调违约和紧急购电费用，并判断增加各次预报是否具有经济价值。

\textbf{问题四：}在实时波动电价下重新计算问题二和问题三，比较因果价格预测与完美日内价格信息，并分析价格波动对储能调度和费用的影响。
